% Generated draft: requires mathematical and source review.
\documentclass[11pt]{article}
\usepackage[margin=25mm]{geometry}
\usepackage{fontspec}
\setmainfont{Latin Modern Roman}
\usepackage{amsmath,amssymb,mathtools,unicode-math}
\setmathfont{Latin Modern Math}
\usepackage{enumitem,needspace,xurl,hyperref}
\setlength{\parindent}{0pt}
\setlength{\parskip}{5pt}
\setlength{\emergencystretch}{3em}
\newcommand{\N}{\mathbb{N}}
\newcommand{\Z}{\mathbb{Z}}
\newcommand{\Q}{\mathbb{Q}}
\newcommand{\R}{\mathbb{R}}
\newcommand{\C}{\mathbb{C}}
\newcommand{\F}{\mathbb{F}}
\newcommand{\A}{\mathbb{A}}
\newcommand{\PP}{\mathbb{P}}
\newcommand{\E}{\mathbb{E}}
\newcommand{\eps}{\varepsilon}
\newcommand{\dd}{\,\mathrm{d}}
\begin{document}
\section*{UnsolvedMath ID: 30005569}
\textbf{Source problem code:} catalogue entry
\subsection{Definitions and mathematical statement}
The supplied question concerns the product
\[
Y=S_K(\operatorname{GSp}_4)\times X_0(N),
\]
where \(S_K(\operatorname{GSp}_4)\) denotes the Shimura variety at level \(K\) appearing in the source, and \(X_0(N)\) denotes the modular curve with \(\Gamma_0(N)\)-level structure. The record does not specify the admissible levels, the base field, the cohomological realization, or the precise meaning of the endoscopic contribution.

\textbf{Intended mathematical question.} Does every endoscopic Tate class in the middle-degree cohomology of \(Y\) arise from an algebraic cycle, including every such class attached to a non-generic member of an endoscopic \(L\)-packet? Here an \(L\)-packet is a collection of representations associated with a common Langlands parameter; genericity refers to the existence of a Whittaker model. The relevant parameters, representations, and local or global genericity conditions are not specified in the supplied record.

A conditional formulation of the cycle-class target is as follows. Suppose the intended data specify a number field \(k\), a prime \(\ell\), a smooth model \(Y/k\) of pure even dimension \(2r\), and ordinary \(\ell\)-adic cohomology. Write \(\bar k\) for an algebraic closure and \(G_{k'}=\operatorname{Gal}(\bar k/k')\) for each finite extension \(k'/k\) in \(\bar k\). Define the space of middle-degree Tate classes by
\[
T^r_\ell(Y)=
\bigcup_{\substack{k\subseteq k'\subseteq\bar k\\[k':k]<\infty}}
H^{2r}_{\mathrm{\acute et}}(Y_{\bar k},\Q_\ell(r))^{G_{k'}}.
\]
Let \(T^r_{\ell,\mathrm{endo}}(Y)\subseteq T^r_\ell(Y)\) denote the endoscopic classes intended by the source. If \(\mathrm{CH}^r(Y_{\bar k})\) is the group of codimension-\(r\) algebraic cycles modulo rational equivalence, the requested algebraicity assertion in this realization is
\[
\forall\alpha\in T^r_{\ell,\mathrm{endo}}(Y),\qquad
\alpha\in\operatorname{im}\!\left(
\mathrm{cl}_\ell:\mathrm{CH}^r(Y_{\bar k})\otimes_{\Z}\Q_\ell
\longrightarrow H^{2r}_{\mathrm{\acute et}}(Y_{\bar k},\Q_\ell(r))
\right).
\]
This quantifier includes the designated classes associated with non-generic packet members. The displayed formulation is conditional: the supplied record does not establish that this model, parity, coefficient convention, or cohomology theory is the one intended. In particular, it does not specify whether a compactification, intersection cohomology, or another realization is required, and it gives no definition of \(T^r_{\ell,\mathrm{endo}}(Y)\).

The supplied dated assessment states that a natural algebraic cycle realizes exactly the generic endoscopic Tate classes and that Hodge-cycle realizations are constructed in the non-generic case, while algebraicity of the non-generic classes remains open. These are supplied literature qualifications, not independently verified conclusions; the reported Hodge-cycle realizations do not themselves assert algebraicity.
\subsection{Short English statement}
The question asks whether all the specified endoscopic Tate classes on \(S_K(\operatorname{GSp}_4)\times X_0(N)\) are algebraic, including those associated with non-generic packet members. The supplied assessment reports the generic case as established and leaves algebraicity in the non-generic case open. Missing geometric and representation-theoretic hypotheses prevent a fully specified unconditional formulation.
\subsection{Sources}
\begin{itemize}
\item Supplied local UnsolvedMath registry record, permanent ID \(30005569\), title \emph{Algebraicity of Endoscopic Tate Classes}, exact source code \texttt{catalogue entry}, in the Oberwolfach Reports Open Problems collection. The record supplies proposed year 2023 and names no proposer.
\item \emph{Algebraische Zahlentheorie} (2023), Oberwolfach Reports, pp.\ 1543--1606, DOI 10.4171/owr/2023/28. Supplied report citation and source URL: \url{https://doi.org/10.4171/owr/2023/28}. The background describes Oberwolfach Reports as records of mathematical workshops and their open-problem discussions.
\item Sweeting, \emph{Tate classes and endoscopy for GSp(4) over totally real fields}. Supplied literature source: \url{https://arxiv.org/abs/2211.10838}.
\item Sweeting, \emph{Tate classes and endoscopy for GSp4}. Supplied literature source: \url{https://doi.org/10.4171/owr/2023/28}.
\item Literature assessment dated exactly \texttt{2026-08-22}: the supplied curation labels the problem partially solved, reports algebraic realization of the generic endoscopic Tate classes and Hodge-cycle realizations in the non-generic case, and states that algebraicity of the latter remains open. This is supplied dated curation, not independently verified. The background explicitly requests independent verification of time-sensitive status and cited claims.
\item The registry labels the statement \texttt{corrected\_verified} and reports checking it against the official Oberwolfach report. This verification claim is supplied metadata. The retained original statement asks when unspecified preceding classes arise from algebraic cycles, as predicted by the Tate conjecture; its preceding mathematical context is absent from this record.
\end{itemize}
\end{document}
